\documentclass[letterpaper,10pt]{article}

\usepackage{amsmath}
\usepackage{amssymb}
\usepackage{bm}
\usepackage{graphicx}
\usepackage{booktabs}
\usepackage{hyperref}
\usepackage[margin=0.75in]{geometry}

\newcommand{\R}{\mathbb{R}}
\newcommand{\posv}{\bm{p}}
\newcommand{\velv}{\bm{v}}
\newcommand{\accv}{\bm{a}}
\newcommand{\linmom}{\bm{\ell}}
\newcommand{\angmom}{\bm{h}}
\newcommand{\omegav}{\bm{\omega}}

\title{\LARGE \bf
Unicycle-Pendulum and Dual-Flywheel Reference Generation for Single-Stance WLH Balance Control
}

\author{Junyoung Kim}

\begin{document}
\maketitle

\begin{abstract}
This note describes a reference generator for a wheeled-legged humanoid that maintains a single rolling contact and is controlled through a unicycle-pendulum model with two reaction flywheels. Given twist command using PD control law, from \((0,0,0)\) to \((1,1,\pi/2)\) into references for the wheel support point, center of mass (CoM), centroidal linear momentum, centroidal angular momentum, wheel speed, and roll/yaw flywheel speeds. The resulting trajectory is intended to be used as \(\mathbf{X}^{\mathrm{ref}}\) for a centroidal MPC or as motion/momentum references for a whole-body controller.
\end{abstract}

\section{Modeling Idea}
The target behavior is a single-contact wheeled-legged humanoid motion. The support contact is treated as a rolling unicycle contact, while the upper body behaves like a pendulum above that contact. The drive wheel controls forward and backward progression. Two flywheels are used as angular-momentum actuators: one mainly regulates roll momentum and the other regulates yaw momentum. Pitch is represented as a body lean reference caused by longitudinal acceleration, although an additional pitch-axis flywheel could be added with the same construction.

Let the planar unicycle state be
\begin{equation}
\bm{x}_u =
\begin{bmatrix}
x_w & y_w & \psi & v
\end{bmatrix}^{\top},
\end{equation}
where \((x_w,y_w)\) is the virtual wheel-contact point and \(\psi\) is the yaw angle. 

\end{document}